\textbf{Phần 1. Trắc nghiệm lựa chọn}
\setcounterpageref{ex}{0}

\begin{ex} %[3H3N19]
Hình có 3 đỉnh và 3 cạnh là
\choice
{Hình tứ giác}
{\True Hình tam giác}
{Hình chữ nhật}
{Hình vuông}
\loigiai{
Hình tam giác có 3 đỉnh và 3 cạnh.\\
\textbf{Đáp án của câu hỏi là B}
}
\end{ex}

\begin{ex} %[3H3N19]
Hình có 4 đỉnh và 4 cạnh là
\choice
{Hình tam giác}
{Hình tròn}
{\True Hình tứ giác}
{Hình ngũ giác}
\loigiai{
Hình tứ giác có 4 đỉnh và 4 cạnh.\\
\textbf{Đáp án của câu hỏi là C}
}
\end{ex}

\begin{ex} %[3H3H19]
Hình chữ nhật có
\choice
{4 cạnh bằng nhau}
{\True 4 góc vuông, 2 cạnh dài bằng nhau và 2 cạnh ngắn bằng nhau}
{Chỉ có 2 góc vuông}
{3 góc vuông}
\loigiai{
Hình chữ nhật có 4 góc vuông, 2 cạnh dài bằng nhau và 2 cạnh ngắn bằng nhau.\\
\textbf{Đáp án của câu hỏi là B}
}
\end{ex}

\begin{ex} %[3H3H19]
Hình vuông là hình tứ giác có
\choice
{4 cạnh bằng nhau nhưng không có góc vuông}
{2 góc vuông và 2 cạnh bằng nhau}
{\True 4 góc vuông và 4 cạnh có độ dài bằng nhau}
{Chỉ có 3 cạnh bằng nhau}
\loigiai{
Hình vuông có 4 góc vuông và 4 cạnh bằng nhau.\\
\textbf{Đáp án của câu hỏi là C}
}
\end{ex}

\begin{ex} %[3H3H19]
Chu vi của hình tam giác có 3 cạnh $a, b, c$ là
\choice
{\True $a + b + c$}
{$a \times b \times c$}
{$(a + b) \times c$}
{$(a + b + c) \times 2$}
\loigiai{
Chu vi tam giác là tổng độ dài 3 cạnh: $a + b + c$.\\
\textbf{Đáp án của câu hỏi là A}
}
\end{ex}

\begin{ex} %[3H3H19]
Chu vi hình vuông có cạnh dài $a$ là
\choice
{$a + 4$}
{\True $a \times 4$}
{$a \times a$}
{$(a + a) \times 2$}
\loigiai{
Chu vi hình vuông là $a \times 4$.\\
\textbf{Đáp án của câu hỏi là B}
}
\end{ex}

\begin{ex} %[3H3V19]
Một hình vuông có cạnh dài $8\text{ cm}$. Chu vi hình vuông đó là
\choice
{$24\text{ cm}$}
{$36\text{ cm}$}
{\True $32\text{ cm}$}
{$16\text{ cm}$}
\loigiai{
$8 \times 4 = 32\text{ cm}$.\\
\textbf{Đáp án của câu hỏi là C}
}
\end{ex}

\begin{ex} %[3H3V19]
Một hình chữ nhật có chiều dài $12\text{ cm}$, chiều rộng $8\text{ cm}$. Chu vi hình chữ nhật đó là
\choice
{$20\text{ cm}$}
{$30\text{ cm}$}
{$96\text{ cm}$}
{\True $40\text{ cm}$}
\loigiai{
$(12 + 8) \times 2 = 40\text{ cm}$.\\
\textbf{Đáp án của câu hỏi là D}
}
\end{ex}

\begin{ex} %[3H3V19]
Một hình tam giác có 3 cạnh bằng nhau và có chu vi là $27\text{ cm}$. Độ dài mỗi cạnh của hình tam giác là
\choice
{\True $9\text{ cm}$}
{$8\text{ cm}$}
{$7\text{ cm}$}
{$6\text{ cm}$}
\loigiai{
$27 : 3 = 9\text{ cm}$.\\
\textbf{Đáp án của câu hỏi là A}
}
\end{ex}

\begin{ex} %[3H3V19]
Chu vi hình vuông là $28\text{ dm}$. Độ dài cạnh hình vuông đó là
\choice
{$6\text{ dm}$}
{$8\text{ dm}$}
{$14\text{ dm}$}
{\True $7\text{ dm}$}
\loigiai{
$28 : 4 = 7\text{ dm}$.\\
\textbf{Đáp án của câu hỏi là D}
}
\end{ex}

\begin{ex} %[3H3V19]
Một hình chữ nhật có nửa chu vi là $18\text{ cm}$, chiều dài $11\text{ cm}$. Chiều rộng của hình chữ nhật là
\choice
{\True $7\text{ cm}$}
{$6\text{ cm}$}
{$8\text{ cm}$}
{$9\text{ cm}$}
\loigiai{
Chiều rộng: $18 - 11 = 7\text{ cm}$.\\
\textbf{Đáp án của câu hỏi là A}
}
\end{ex}

\begin{ex} %[3H3C19]
Một mảnh đất hình chữ nhật có chiều rộng $6\text{ m}$, chiều dài gấp 3 lần chiều rộng. Chu vi mảnh đất đó là
\choice
{$36\text{ m}$}
{$42\text{ m}$}
{$54\text{ m}$}
{\True $48\text{ m}$}
\loigiai{
Chiều dài: $6 \times 3 = 18\text{ m}$.\\
Chu vi: $(18 + 6) \times 2 = 48\text{ m}$.\\
\textbf{Đáp án của câu hỏi là D}
}
\end{ex}

\textbf{Phần 2. Trắc nghiệm đúng sai}
\setcounterpageref{ex}{0}

\begin{ex} %[3H3H19]
Xét tính đúng sai của các khẳng định về hình học:
\choiceTF
{\True Hình vuông cũng là một hình chữ nhật đặc biệt có 4 cạnh bằng nhau.}
{\True Hình chữ nhật có 4 góc vuông.}
{Hình tam giác có 4 đỉnh và 3 cạnh.}
{\True Hình tam giác có 3 đỉnh và 3 cạnh.}
\loigiai{
\begin{itemize}
    \item Hình vuông là hình chữ nhật đặc biệt: Đúng.
    \item Hình chữ nhật có 4 góc vuông: Đúng.
    \item Tam giác có 3 đỉnh (không phải 4). Ý 3 Sai, Ý 4 Đúng.
\end{itemize}
\begin{itemize}
    \item a) \textbf{Đúng}.
    \item b) \textbf{Đúng}.
    \item c) \textbf{Sai}.
    \item d) \textbf{Đúng}.
\end{itemize}
}
\end{ex}

\begin{ex} %[3H3V19]
Cho hình chữ nhật $MNPQ$ có $MN = 10\text{ cm}, NP = 5\text{ cm}$. Xét tính đúng sai:
\choiceTF
{\True Cạnh $PQ = 10\text{ cm}$.}
{\True Cạnh $MQ = 5\text{ cm}$.}
{Chu vi hình chữ nhật $MNPQ$ là $15\text{ cm}$.}
{\True Chu vi hình chữ nhật $MNPQ$ là $30\text{ cm}$.}
\loigiai{
\begin{itemize}
    \item $PQ = MN = 10\text{ cm}$. Đúng.
    \item $MQ = NP = 5\text{ cm}$. Đúng.
    \item Chu vi: $(10 + 5) \times 2 = 30\text{ cm} \neq 15\text{ cm}$. Ý 3 Sai, Ý 4 Đúng.
\end{itemize}
\begin{itemize}
    \item a) \textbf{Đúng}.
    \item b) \textbf{Đúng}.
    \item c) \textbf{Sai}.
    \item d) \textbf{Đúng}.
\end{itemize}
}
\end{ex}

\textbf{Phần 3. Trả lời ngắn}
\setcounterpageref{ex}{0}

\begin{ex} %[3H3H19]
Tính chu vi hình vuông có cạnh dài $9\text{ cm}$.
\shortans{36}
\loigiai{
\textbf{Đáp số:} $36$\\
$9 \times 4 = 36\text{ cm}$.
}
\end{ex}

\begin{ex} %[3H3H19]
Tính chu vi hình chữ nhật có chiều dài $15\text{ cm}$, chiều rộng $5\text{ cm}$.
\shortans{40}
\loigiai{
\textbf{Đáp số:} $40$\\
$(15 + 5) \times 2 = 40\text{ cm}$.
}
\end{ex}

\begin{ex} %[3H3V19]
Tính chu vi tam giác có 3 cạnh dài $7\text{ cm}, 8\text{ cm}, 9\text{ cm}$.
\shortans{24}
\loigiai{
\textbf{Đáp số:} $24$\\
$7 + 8 + 9 = 24\text{ cm}$.
}
\end{ex}

\begin{ex} %[3H3V19]
Một hình vuông có chu vi $32\text{ dm}$. Tính độ dài cạnh hình vuông theo đề-xi-mét.
\shortans{8}
\loigiai{
\textbf{Đáp số:} $8$\\
$32 : 4 = 8\text{ dm}$.
}
\end{ex}

\begin{ex} %[3H3V19]
Một hình chữ nhật có chu vi $24\text{ cm}$, chiều dài là $8\text{ cm}$. Tính chiều rộng của hình chữ nhật theo xăng-ti-mét.
\shortans{4}
\loigiai{
\textbf{Đáp số:} $4$\\
Nửa chu vi: $24 : 2 = 12\text{ cm}$. Chiều rộng: $12 - 8 = 4\text{ cm}$.
}
\end{ex}

\begin{ex} %[3H3C19]
Ghép 2 hình vuông cạnh $4\text{ cm}$ lại với nhau thành một hình chữ nhật. Tính chu vi hình chữ nhật vừa tạo thành theo xăng-ti-mét.
\shortans{24}
\loigiai{
\textbf{Đáp số:} $24$\\
Chiều dài: $4 + 4 = 8\text{ cm}$, Chiều rộng: $4\text{ cm}$. Chu vi: $(8 + 4) \times 2 = 24\text{ cm}$.
}
\end{ex}
